%% ma: L12-B11-0016
%% lop: 12
%% chuong: 4
%% bai: 11
%% dang: 12.11.12
%% loai: TL
%% muc: TH
%% dap_an: "-"
%% nguon: "(NBV)-10 ĐỀ KIỂM TRA HK2-TOÁN 12-DUNG CHUNG 3 SÁCH.pdf (H01, Đề số 1), Phần 4 Câu 1 (tự luận)"
%% kien_thuc: "Bài 11 (nguyên hàm: tìm vận tốc, quãng đường từ gia tốc)"
%% trang_thai: chinh-thuc
%% ghi_chu: "Câu tự luận. Đáp án gốc 1000/49 ≈ 20,41 m đúng (giải độc lập). Bỏ ảnh nhà khoa học (ảnh trang trí). Đã duyệt 19/10/2026."
\begin{ex}
Một nhà khoa học tự chế tên lửa và phóng tên lửa từ mặt đất với vận tốc ban đầu là $20$ m/s. Giả sử bỏ qua sức cản của gió, tên lửa chỉ chịu tác động của trọng lực. Hỏi độ cao lớn nhất mà tên lửa đạt được là bao nhiêu mét?
\loigiai{Chọn gốc thời gian lúc phóng, chiều dương hướng lên. Gọi $s(t)$, $v(t)$ là độ cao và vận tốc của tên lửa tại thời điểm $t$; ta có $s(0)=0$, $v(0)=20$. Gia tốc là $a(t)=-9{,}8\ \mathrm{m/s^2}$.

Vận tốc là nguyên hàm của gia tốc: $v(t)=\displaystyle\int(-9{,}8)\,\mathrm dt=-9{,}8t+C_1$; từ $v(0)=20$ suy ra $C_1=20$, $v(t)=-9{,}8t+20$.

Độ cao là nguyên hàm của vận tốc: $s(t)=\displaystyle\int(-9{,}8t+20)\,\mathrm dt=-4{,}9t^2+20t+C_2$; từ $s(0)=0$ suy ra $C_2=0$, $s(t)=-4{,}9t^2+20t$.

Tên lửa lên cao nhất khi $v(t)=0$, tức $t=\dfrac{20}{9{,}8}=\dfrac{100}{49}$ (s). Khi đó $s\left(\dfrac{100}{49}\right)=\dfrac{1000}{49}\approx20{,}41$ m.}
\end{ex}
